%---------------------------------------------------------------------
%
%                          acronimos.tex
%
%---------------------------------------------------------------------
%
% Lista de acrónimos usados en la memoria. Se mantiene a mano, en el
% mismo bloque de índices que las listas de figuras y tablas.
%
%---------------------------------------------------------------------

\ifpdf
   \pdfbookmark[1]{Lista de acrónimos}{lista de acronimos}
\fi
\cabeceraEspecial{Lista de acr\'onimos}
\chapter*{Lista de acrónimos}

\renewcommand{\arraystretch}{1.1}
\begin{tabular}{@{}p{2.3cm}p{11.4cm}@{}}
ADE & Administración y Dirección de Empresas. \\
API & \textit{Application Programming Interface}, interfaz de programación de aplicaciones. \\
BEDCA & Base de Datos Española de Composición de Alimentos. \\
BMR & \textit{Basal Metabolic Rate}, metabolismo basal. \\
CP & \textit{Constraint Programming}, programación con restricciones. \\
CP-SAT & Solver de programación con restricciones de OR-Tools, que combina propagación de restricciones con técnicas SAT. \\
CPU & \textit{Central Processing Unit}, unidad central de procesamiento. \\
CSS & \textit{Cascading Style Sheets}, hojas de estilo de la web. \\
HIPAA & \textit{Health Insurance Portability and Accountability Act}, ley estadounidense de protección de datos sanitarios. \\
IA & Inteligencia artificial. \\
JSON & \textit{JavaScript Object Notation}, formato de intercambio de datos. \\
JSONB & Tipo de columna de PostgreSQL que guarda documentos JSON en formato binario. \\
LLM & \textit{Large Language Model}, modelo de lenguaje de gran tamaño. \\
MCP & \textit{Model Context Protocol}, protocolo que conecta un asistente de IA con herramientas externas. \\
PDF & \textit{Portable Document Format}. \\
RF & Requisito funcional. \\
RGPD & Reglamento General de Protección de Datos. \\
RLS & \textit{Row Level Security}, seguridad a nivel de fila. \\
RNF & Requisito no funcional. \\
SAT & \textit{Boolean Satisfiability}, problema de satisfacibilidad booleana. \\
SQL & \textit{Structured Query Language}, lenguaje de consulta de bases de datos relacionales. \\
TFG & Trabajo de Fin de Grado. \\
URL & \textit{Uniform Resource Locator}, dirección de un recurso en la web. \\
USDA & \textit{United States Department of Agriculture}, referencia internacional en composición de alimentos. \\
\end{tabular}
\renewcommand{\arraystretch}{1}

\newpage
